% Options for packages loaded elsewhere
% Options for packages loaded elsewhere
\PassOptionsToPackage{unicode}{hyperref}
\PassOptionsToPackage{hyphens}{url}
%
\documentclass[
  ignorenonframetext,
  aspectratio=169,
  brazilian,
  t,
  10pt]{beamer}
\newif\ifbibliography
\usepackage{pgfpages}
\setbeamertemplate{caption}[numbered]
\setbeamertemplate{caption label separator}{: }
\setbeamercolor{caption name}{fg=normal text.fg}
\beamertemplatenavigationsymbolsempty
% remove section numbering
\setbeamertemplate{part page}{
  \centering
  \begin{beamercolorbox}[sep=16pt,center]{part title}
    \usebeamerfont{part title}\insertpart\par
  \end{beamercolorbox}
}
\setbeamertemplate{section page}{
  \centering
  \begin{beamercolorbox}[sep=12pt,center]{section title}
    \usebeamerfont{section title}\insertsection\par
  \end{beamercolorbox}
}
\setbeamertemplate{subsection page}{
  \centering
  \begin{beamercolorbox}[sep=8pt,center]{subsection title}
    \usebeamerfont{subsection title}\insertsubsection\par
  \end{beamercolorbox}
}
% Prevent slide breaks in the middle of a paragraph
\widowpenalties 1 10000
\raggedbottom
\AtBeginPart{
  \frame{\partpage}
}
\AtBeginSection{
  \ifbibliography
  \else
    \frame{\sectionpage}
  \fi
}
\AtBeginSubsection{
  \frame{\subsectionpage}
}
\usepackage{iftex}
\ifPDFTeX
  \usepackage[T1]{fontenc}
  \usepackage[utf8]{inputenc}
  \usepackage{textcomp} % provide euro and other symbols
\else % if luatex or xetex
  \usepackage{unicode-math} % this also loads fontspec
  \defaultfontfeatures{Scale=MatchLowercase}
  \defaultfontfeatures[\rmfamily]{Ligatures=TeX,Scale=1}
\fi
\usepackage{lmodern}

\usetheme[]{Madrid}
\usecolortheme[]{default}
\usefonttheme[]{professionalfonts}
\ifPDFTeX\else
  % xetex/luatex font selection
\fi
% Use upquote if available, for straight quotes in verbatim environments
\IfFileExists{upquote.sty}{\usepackage{upquote}}{}
\IfFileExists{microtype.sty}{% use microtype if available
  \usepackage[]{microtype}
  \UseMicrotypeSet[protrusion]{basicmath} % disable protrusion for tt fonts
}{}
\makeatletter
\@ifundefined{KOMAClassName}{% if non-KOMA class
  \IfFileExists{parskip.sty}{%
    \usepackage{parskip}
  }{% else
    \setlength{\parindent}{0pt}
    \setlength{\parskip}{6pt plus 2pt minus 1pt}}
}{% if KOMA class
  \KOMAoptions{parskip=half}}
\makeatother


\usepackage{longtable,booktabs,array}
\newcounter{none} % for unnumbered tables
\usepackage{calc} % for calculating minipage widths
\usepackage{caption}
% Make caption package work with longtable
\makeatletter
\def\fnum@table{\tablename~\thetable}
\makeatother
\usepackage{graphicx}
\makeatletter
\newsavebox\pandoc@box
\newcommand*\pandocbounded[1]{% scales image to fit in text height/width
  \sbox\pandoc@box{#1}%
  \Gscale@div\@tempa{\textheight}{\dimexpr\ht\pandoc@box+\dp\pandoc@box\relax}%
  \Gscale@div\@tempb{\linewidth}{\wd\pandoc@box}%
  \ifdim\@tempb\p@<\@tempa\p@\let\@tempa\@tempb\fi% select the smaller of both
  \ifdim\@tempa\p@<\p@\scalebox{\@tempa}{\usebox\pandoc@box}%
  \else\usebox{\pandoc@box}%
  \fi%
}
% Set default figure placement to htbp
\def\fps@figure{htbp}
\makeatother



\ifLuaTeX
\usepackage[bidi=basic,shorthands=off]{babel}
\else
\usepackage[bidi=default,shorthands=off]{babel}
\fi
\ifLuaTeX
  \usepackage{selnolig} % disable illegal ligatures
\fi


\setlength{\emergencystretch}{3em} % prevent overfull lines

\providecommand{\tightlist}{%
  \setlength{\itemsep}{0pt}\setlength{\parskip}{0pt}}



 


\useoutertheme{miniframes}
\useinnertheme{rounded}
\usepackage{booktabs}
\usepackage{tabularx}
\usepackage[most]{tcolorbox}
\author[T. Mançano | de Kadt \& Grzymala-Busse]{Daniel de Kadt \& Anna Grzymala-Busse (2025)}
\AtBeginDocument{%
  \institute[PPGCP-USP]{Universidade de São Paulo --- PPGCP \\ Apresentação: Tales Mançano}%
}
\makeatletter
\@ifpackageloaded{caption}{}{\usepackage{caption}}
\AtBeginDocument{%
\ifdefined\contentsname
  \renewcommand*\contentsname{Índice}
\else
  \newcommand\contentsname{Índice}
\fi
\ifdefined\listfigurename
  \renewcommand*\listfigurename{Lista de Figuras}
\else
  \newcommand\listfigurename{Lista de Figuras}
\fi
\ifdefined\listtablename
  \renewcommand*\listtablename{Lista de Tabelas}
\else
  \newcommand\listtablename{Lista de Tabelas}
\fi
\ifdefined\figurename
  \renewcommand*\figurename{Figura}
\else
  \newcommand\figurename{Figura}
\fi
\ifdefined\tablename
  \renewcommand*\tablename{Tabela}
\else
  \newcommand\tablename{Tabela}
\fi
}
\@ifpackageloaded{float}{}{\usepackage{float}}
\floatstyle{ruled}
\@ifundefined{c@chapter}{\newfloat{codelisting}{h}{lop}}{\newfloat{codelisting}{h}{lop}[chapter]}
\floatname{codelisting}{Listagem}
\newcommand*\listoflistings{\listof{codelisting}{Lista de Listagens}}
\makeatother
\makeatletter
\makeatother
\makeatletter
\@ifpackageloaded{caption}{}{\usepackage{caption}}
\@ifpackageloaded{subcaption}{}{\usepackage{subcaption}}
\makeatother

\usepackage{bookmark}
\IfFileExists{xurl.sty}{\usepackage{xurl}}{} % add URL line breaks if available
\urlstyle{same}
\hypersetup{
  pdftitle={Good Description},
  pdfauthor={Tales Mançano},
  pdflang={pt-BR},
  hidelinks,
  pdfcreator={LaTeX via pandoc}}


\title{\texorpdfstring{Good Description}{Good Description}}
\subtitle{\texorpdfstring{Uma Teoria da Descrição na Ciência Social
Empírica}{Uma Teoria da Descrição na Ciência Social Empírica}}
\author{Daniel de Kadt \& Anna Grzymala-Busse (2025)}
\date{2026-01-01}
\institute{Fichamento e Apresentação: Tales Mançano --- PPGCP-USP}

\begin{document}
\frame{\titlepage}


\section{O Paradoxo e o Propósito da
Descrição}\label{o-paradoxo-e-o-propuxf3sito-da-descriuxe7uxe3o}

\begin{frame}{O Paradoxo da Descrição nas Ciências Sociais}
\protect\phantomsection\label{o-paradoxo-da-descriuxe7uxe3o-nas-ciuxeancias-sociais}
\begin{block}{A Tensão Disciplinar}
\protect\phantomsection\label{a-tensuxe3o-disciplinar}
\begin{itemize}
\tightlist
\item
  \textbf{Desprezo formal}: A descrição é frequentemente derribada ou
  tratada como modo ``inferior'' de investigação em comparação à
  inferência causal (Gerring 2012; Kreuzer 2019; Holmes et al.~2024).
\item
  \textbf{Prática empírica dominante}: Artigos e livros
  predominantemente descritivos continuam sendo amplamente publicados
  nos principais periódicos (\emph{APSR}, \emph{AJPS}, \emph{World
  Politics}) e premiados pela disciplina (Przeworski \& Limongi 1997;
  Beissinger 2022; Daly 2022; Stasavage 2020).
\item
  \textbf{A citação de Grimmer (2015, p.~80)}: \emph{``Muito do trabalho
  empírico dos cientistas políticos e das teorias que eles constroem é
  um produto direto da descrição.''}
\end{itemize}
\end{block}

\begin{block}{A Origem do Paradoxo}
\protect\phantomsection\label{a-origem-do-paradoxo}
\begin{itemize}
\tightlist
\item
  A ciência social carece de um \textbf{framework avaliativo coerente}
  para a descrição.
\item
  Enquanto a inferência causal possui manuais, \emph{checklists} e
  critérios formais maduros, a descrição padece de incerteza sobre o que
  constitui uma \textbf{boa pergunta} e uma \textbf{boa análise}.
\end{itemize}
\end{block}
\end{frame}

\begin{frame}{O Objetivo do Artigo e a Tese Central}
\protect\phantomsection\label{o-objetivo-do-artigo-e-a-tese-central}
\begin{block}{Objetivos Principais}
\protect\phantomsection\label{objetivos-principais}
\begin{enumerate}
\tightlist
\item
  \textbf{Articular o propósito científico da descrição}: O que torna
  uma pergunta descritiva cientificamente relevante?
\item
  \textbf{Definir as características ideais da análise descritiva}:
  Quais critérios permitem distinguir a ``boa'' da ``mera'' descrição?
\end{enumerate}

\begin{block}{Tese Central do Artigo}
Descrição e inferência causal \textbf{não são modos epistemologicamente distintos} de investigação. Ambas contribuem para o mesmo objetivo científico fundamental: o desenvolvimento incremental, o refinamento e a avaliação de teorias e explicações do mundo social.
\end{block}
\end{block}

\begin{block}{Unidade de Propósito Epistêmico}
\protect\phantomsection\label{unidade-de-propuxf3sito-epistuxeamico}
\begin{itemize}
\tightlist
\item
  A descrição de excelência é \textbf{invariavelmente orientada pela
  teoria} (\emph{theory-led}).
\item
  Uma boa pergunta descritiva busca descobrir fatos que exigem
  explicação, construir novas teorias ou avaliar implicações de teorias
  existentes.
\end{itemize}
\end{block}
\end{frame}

\section{Fundamentação Teórica e a ``Escada'' de
Perguntas}\label{fundamentauxe7uxe3o-teuxf3rica-e-a-escada-de-perguntas}

\begin{frame}{O Critério de Holland: Manipulabilidade}
\protect\phantomsection\label{o-crituxe9rio-de-holland-manipulabilidade}
\begin{block}{A Crítica às Definições Baseadas em Linguagem Natural}
\protect\phantomsection\label{a-cruxedtica-uxe0s-definiuxe7uxf5es-baseadas-em-linguagem-natural}
\begin{itemize}
\tightlist
\item
  Definições tradicionais diferenciam descrição (``quem, o quê, onde,
  quando'') de causação (``por quê'') (Gerring 2012; Holmes et
  al.~2024).
\item
  \textbf{Problema de imprecisão}: Perguntas como \emph{``O que
  aconteceu quando a intervenção X ocorreu?''} usam a linguagem de ``o
  quê'', mas são puramente causais.
\end{itemize}
\end{block}

\begin{block}{A Solução Formal (Holland 1986)}
\protect\phantomsection\label{a-soluuxe7uxe3o-formal-holland-1986}
\begin{itemize}
\tightlist
\item
  \textbf{Critério de Manipulabilidade}: Uma pergunta é \textbf{causal}
  se exige a manipulabilidade de pelo menos uma variável para ser
  respondida. Se a manipulabilidade \textbf{não} é necessária, a
  pergunta é \textbf{descritiva}.
\end{itemize}

\begin{exampleblock}{Exemplo Substantivo: Raça vs. Discriminação Racial (Sen \& Wasow 2016)}
\begin{itemize}
  \item \textbf{Pergunta Descritiva}: \textit{Pessoas de diferentes grupos raciais ganham salários diferentes?} (Raça não é manipulável; pergunta sobre o estado real do mundo).
  \item \textbf{Pergunta Causal}: \textit{A discriminação no mercado de trabalho ou uma política pública altera a lacuna salarial?} (Intervenção manipulável; raciocínio contrafactual).
\end{itemize}
\end{exampleblock}
\end{block}
\end{frame}

\begin{frame}{A ``Escada'' de Perguntas de Pesquisa (DAG \(A \to B\))}
\protect\phantomsection\label{a-escada-de-perguntas-de-pesquisa-dag-a-to-b}
\begin{block}{Seis Degraus da Investigação Científica (Figura 1 do Artigo)}
\begin{enumerate}
  \item \textbf{Conceito}: O que são $A$ e $B$? (Estatuto ontológico e empírico).
  \item \textbf{Característica}: Com que frequência ou magnitude $A$ e $B$ ocorrem?
  \item \textbf{Associação}: $A$ e $B$ co-ocorrem na população?
  \item \textbf{Associação Condicional}: $A$ e $B$ co-ocorrem \textit{ceteris paribus}?
  \item \textbf{Efeito Causal}: $A$ causa uma mudança na probabilidade ou nível de $B$?
  \item \textbf{Mecanismo Explicativo}: De que maneira $A$ produz o efeito em $B$?
\end{enumerate}
\end{block}

\begin{block}{Propriedades da Escada}
\protect\phantomsection\label{propriedades-da-escada}
\begin{itemize}
\tightlist
\item
  \textbf{Hierarquia Lógica}: Não se pode responder de forma crível a
  perguntas de degraus inferiores sem solucionar as perguntas dos
  degraus superiores.
\item
  \textbf{Rungs 1 a 3}: Puramente descritivos.
\item
  \textbf{Rung 4 (Associação Condicional)}: Zona liminar ---
  frequentemente usada como substituto euphemístico de causação (Hernán
  2018).
\end{itemize}
\end{block}
\end{frame}

\begin{frame}{Teoria Causal vs.~Pergunta de Pesquisa}
\protect\phantomsection\label{teoria-causal-vs.-pergunta-de-pesquisa}
\begin{block}{Distinção Crucial de Níveis}
\protect\phantomsection\label{distinuxe7uxe3o-crucial-de-nuxedveis}
\begin{itemize}
\tightlist
\item
  \textbf{Teoria Subjacente}: Quase sempre causal (\(A \to B\)).
\item
  \textbf{Pergunta de Pesquisa Alvo}: Pode ser intermediária e
  descritiva (ex.: medir a frequência de \(A\) ou a associação
  condicional de \(A\) e \(B\)).
\end{itemize}
\end{block}

\begin{block}{Evitando o Uso Eufemístico da Linguagem}
\protect\phantomsection\label{evitando-o-uso-eufemuxedstico-da-linguagem}
\begin{itemize}
\tightlist
\item
  Pesquisadores frequentemente usam linguagem descritiva como
  ``eufemismo'' por falta de identificação causal adequada (Hernán 2018;
  Haber et al.~2022).
\item
  \textbf{Diretriz}: Se a evidência disponível só permite responder a
  perguntas descritivas, deve-se formular abertamente \textbf{boas
  perguntas descritivas}, em vez de fingir causalidade fraca.
\end{itemize}
\end{block}
\end{frame}

\section{As Três Metas Científicas da Boa
Descrição}\label{as-truxeas-metas-cientuxedficas-da-boa-descriuxe7uxe3o}

\begin{frame}{Meta 1: Refletir e Medir Teoria}
\protect\phantomsection\label{meta-1-refletir-e-medir-teoria}
\begin{block}{Medição e Validade de Construto}
\protect\phantomsection\label{mediuxe7uxe3o-e-validade-de-construto}
\begin{itemize}
\tightlist
\item
  A medição é a função descritiva mais consensualmente aceita nas
  ciências sociais.
\item
  \textbf{Orientação Teórica Necessária}: Observações não são
  independentes da teoria (Kuhn 1970; Lakatos 1970; Hall 2003).
\item
  Exemplo: Índices de democracia (Przeworski \& Limongi 1997; V-Dem) ou
  polarização urbano-rural (Huijsmans \& Rodden 2024).
\end{itemize}
\end{block}

\begin{block}{Medição de ``Primitivos'' Sociais}
\protect\phantomsection\label{mediuxe7uxe3o-de-primitivos-sociais}
\begin{itemize}
\tightlist
\item
  Conceitos fundamentais (\emph{primitivos}) como riqueza, tipo de
  regime ou desigualdade exigem esforço descritivo intenso.
\item
  \textbf{Validade}: Medidas com baixa validade de construto comprometem
  tanto a descrição quanto a inferência causal subsequente.
\end{itemize}
\end{block}
\end{frame}

\begin{frame}{Meta 2: Construir Teoria e Descobrir Padrões}
\protect\phantomsection\label{meta-2-construir-teoria-e-descobrir-padruxf5es}
\begin{block}{Revelando Anomalias e Quebra-Cabeças}
\protect\phantomsection\label{revelando-anomalias-e-quebra-cabeuxe7as}
\begin{itemize}
\tightlist
\item
  A descrição abrangente permite identificar padrões, \emph{outliers} e
  anomalias que a teoria vigente não previra.
\item
  Anomalias só existem contra o pano de fundo de expectativas teóricas
  prévias.
\end{itemize}
\end{block}

\begin{block}{Iteração Abdutiva / Recursiva}
\protect\phantomsection\label{iterauxe7uxe3o-abdutiva-recursiva}
\begin{itemize}
\tightlist
\item
  Movimento contínuo e transparente entre dados empíricos e modelos
  causais (Geddes 2003; Tavory \& Timmermans 2014; Humphreys \& Jacobs
  2023).
\item
  A observação cuidadosa da realidade impulsiona a formulação de novas
  hipóteses mecanísticas.
\end{itemize}
\end{block}
\end{frame}

\begin{frame}{Meta 3: Avaliar Teoria e Desconfirmar Implicações}
\protect\phantomsection\label{meta-3-avaliar-teoria-e-desconfirmar-implicauxe7uxf5es}
\begin{block}{Testando Implicações Não-Causais de Teorias Causais}
\protect\phantomsection\label{testando-implicauxe7uxf5es-nuxe3o-causais-de-teorias-causais}
\begin{itemize}
\tightlist
\item
  Teorias causais formulam predições sobre correlações e estados do
  mundo observáveis.
\item
  A descrição permite testar se essas implicações de fundo se sustentam
  na realidade.
\end{itemize}

\begin{block}{O Poder Assimétrico da Desconfirmação}
Se uma teoria prediz uma co-ocorrência entre $A$ e $B$ e a análise descritiva demonstra que tal co-ocorrência \textbf{não existe} (ou ocorre em direção oposta), a teoria é imediatamente falsificada ou revelada incompleta.
\end{block}
\end{block}

\begin{block}{Avaliação de Premissas Metodológicas}
\protect\phantomsection\label{avaliauxe7uxe3o-de-premissas-metodoluxf3gicas}
\begin{itemize}
\tightlist
\item
  A descrição histórica e contextual pode invalidar premissas de
  identificação causal (ex.: \emph{spatial RDD} em fronteiras coloniais
  na África; Kocher \& Monteiro 2016; Paine et al.~2025).
\end{itemize}
\end{block}
\end{frame}

\section{As Três Características da Boa Análise
Descritiva}\label{as-truxeas-caracteruxedsticas-da-boa-anuxe1lise-descritiva}

\begin{frame}{Os Três Critérios de Qualidade Analítica}
\protect\phantomsection\label{os-truxeas-crituxe9rios-de-qualidade-analuxedtica}
\begin{block}{Triad Avaliativa da Boa Descrição (Seção 3 do Artigo)}
Para ser cientificamente valiosa, qualquer análise descritiva deve cumprir três critérios:
\begin{enumerate}
  \item \textbf{Clareza} (\textit{Clear}): Conexão transparente entre teoria, pergunta e método.
  \item \textbf{Comparabilidade} (\textit{Comparable}): Sensibilidade ao contexto e significado dos dados.
  \item \textbf{Completude} (\textit{Complete}): Abrangência populacional e analítica.
\end{enumerate}
\end{block}
\end{frame}

\begin{frame}{Clareza e Comparabilidade no Contexto}
\protect\phantomsection\label{clareza-e-comparabilidade-no-contexto}
\begin{block}{Clareza Analítica}
\protect\phantomsection\label{clareza-analuxedtica}
\begin{itemize}
\tightlist
\item
  Alinhamento explícito: Por que descrever este fenômeno? Qual é a
  pergunta exata? Como o método escolhido responde à pergunta?
\end{itemize}
\end{block}

\begin{block}{Comparabilidade e Sensibilidade Contextual}
\protect\phantomsection\label{comparabilidade-e-sensibilidade-contextual}
\begin{itemize}
\tightlist
\item
  Dados de contextos distintos (tempo, espaço, regime) nem sempre
  possuem a mesma confiabilidade ou o mesmo significado.
\item
  \textbf{Exemplo de Produção de Dados}: Dados de violência contra
  jornalistas na Hungria (CPJ 2025) mostram zero mortes, não por força
  democrática, mas porque 80\% da mídia foi cooptada pelo estado (Carey
  \& Gohdes 2021).
\item
  \textbf{O Trade-off}: Balanço entre generalização padronizada e
  sensibilidade ao contexto local.
\end{itemize}
\end{block}
\end{frame}

\begin{frame}{Completude: População, Escopo e Profundidade}
\protect\phantomsection\label{completude-populauxe7uxe3o-escopo-e-profundidade}
\begin{block}{Três Dimensões da Completude}
\protect\phantomsection\label{truxeas-dimensuxf5es-da-completude}
\begin{enumerate}
\tightlist
\item
  \textbf{Completude de Perguntas}: Responder a tantas perguntas
  derivadas da teoria quanto possível.
\item
  \textbf{Completude Populacional}: Especificação rigorosa da população
  inferencial e amostragem adequada.
\item
  \textbf{Trade-off Amplitude vs.~Profundidade}: Estudo profundo de caso
  único ao longo do tempo vs.~estudo amplo de múltiplos casos.
\end{enumerate}
\end{block}
\end{frame}

\section{Mapeamento Pergunta-Método e o Uso de
Regressão}\label{mapeamento-pergunta-muxe9todo-e-o-uso-de-regressuxe3o}

\begin{frame}{Regressão Multivariada como ``Descrição''?}
\protect\phantomsection\label{regressuxe3o-multivariada-como-descriuxe7uxe3o}
\begin{block}{A Tensão Metodológica}
\protect\phantomsection\label{a-tensuxe3o-metodoluxf3gica}
\begin{itemize}
\tightlist
\item
  Pesquisadores frequentemente usam regressões multivariadas com a
  intenção declarada de ``apenas descrever os dados''.
\item
  \textbf{O Dilema}: A regressão por definição realiza projeções em
  espaços contrafactuais e extrapolações fora do suporte dos dados.
\end{itemize}
\end{block}

\begin{block}{A Advertência de Stanley Lieberson (1987, pp.~213--214)}
\protect\phantomsection\label{a-advertuxeancia-de-stanley-lieberson-1987-pp.-213214}
\begin{quote}
"A variável de controle é um dispositivo descritivo perfeitamente apropriado; o problema ocorre quando variáveis de controle são vistas como um procedimento analítico. É uma coisa perguntar: qual é a realização ocupacional de negros e brancos de um dado nível educacional? É outra perguntar: controlando a educação, qual é a influência da raça na realização?"
\end{quote}
\end{block}
\end{frame}

\begin{frame}{Quando Condicionar Faz Sentido}
\protect\phantomsection\label{quando-condicionar-faz-sentido}
\begin{block}{Casos Legítimos de Condicionamento}
\protect\phantomsection\label{casos-leguxedtimos-de-condicionamento}
\begin{enumerate}
\tightlist
\item
  \textbf{Reflexo de Teoria Condicional}: Quando a própria teoria prediz
  que a associação é condicional a \(Z\).
\item
  \textbf{Evitando Conclusões Enganosas (Paradoxo de Simpson 1951)}:
  Quando o sinal global da associação se inverte nos subgrupos.
\item
  \textbf{Desenhos Cavais Explicitos}: Seleção sobre observáveis /
  \emph{matching}.
\end{enumerate}
\end{block}

\begin{block}{Diretriz Prática}
\protect\phantomsection\label{diretriz-pruxe1tica}
\begin{itemize}
\tightlist
\item
  Se o objetivo é puramente descrever os dados, ferramentas
  não-paramétricas, estatísticas descritivas e estratificação direta
  costumam ser superiores à regressão multivariada.
\end{itemize}
\end{block}
\end{frame}

\section{Boa Descrição na
Prática}\label{boa-descriuxe7uxe3o-na-pruxe1tica}

\begin{frame}{Diversidade de Contextos e Explicações Monocausais}
\protect\phantomsection\label{diversidade-de-contextos-e-explicauxe7uxf5es-monocausais}
\begin{block}{Testando Fatores Confundidores e Escopo}
\protect\phantomsection\label{testando-fatores-confundidores-e-escopo}
\begin{itemize}
\tightlist
\item
  Comparar o mesmo fenômeno em contextos diversos permite desqualificar
  explicações monocausais simplistas.
\end{itemize}

\begin{exampleblock}{Exemplos Empíricos da Literatura}
\begin{itemize}
  \item \textbf{Smartphones e Saúde Mental de Adolescentes} (Haidt 2024 vs. Sacco et al. 2024): Uso de celular é alto nos EUA e na Europa, mas depressão pré-adolescente é menor na Europa $\implies$ tese monocausal é incompleta.
  \item \textbf{SUVs e Mortes no Trânsito} (Burn-Murdoch 2024): Vendas de SUVs dispararam em vários países ricos, mas mortes só cresceram nos EUA $\implies$ o veículo não é a causa única/independente.
\end{itemize}
\end{exampleblock}
\end{block}
\end{frame}

\begin{frame}{Distribuições Estatísticas e Leis de Potência}
\protect\phantomsection\label{distribuiuxe7uxf5es-estatuxedsticas-e-leis-de-potuxeancia}
\begin{block}{A Informação embutida nas Distribuições}
\protect\phantomsection\label{a-informauxe7uxe3o-embutida-nas-distribuiuxe7uxf5es}
\begin{itemize}
\tightlist
\item
  As ciências sociais assumem frequentemente distribuições normais
  (curva de Gauss).
\item
  Fenômenos reais (guerras, riqueza, tamanho de cidades) seguem
  frequentemente \textbf{Leis de Potência} (\(1/x\), Pareto, Zipf).
\end{itemize}
\end{block}

\begin{block}{Implicações Teóricas das Power Laws}
\protect\phantomsection\label{implicauxe7uxf5es-teuxf3ricas-das-power-laws}
\begin{itemize}
\tightlist
\item
  \textbf{Caudas Longas}: Eventos extremos e raros ocorrem com
  frequência muito maior que a prevista por Gauss.
\item
  \textbf{Natureza Fractal}: Conflitos pequenos são versões fractais de
  grandes guerras.
\item
  \textbf{Interdependência}: Leis de potência resultam da
  \emph{multiplicação} e \emph{corrupção/feedback} de variáveis,
  violando premissas de independência (West 2017; Miller \& Page 2009).
\end{itemize}
\end{block}
\end{frame}

\begin{frame}{Desdobramento no Tempo e no Espaço (\emph{Process
Tracing})}
\protect\phantomsection\label{desdobramento-no-tempo-e-no-espauxe7o-process-tracing}
\begin{block}{Descrição Temporal e Espacial}
\protect\phantomsection\label{descriuxe7uxe3o-temporal-e-espacial}
\begin{itemize}
\tightlist
\item
  Rastrear a difusão, duração, ritmo e sequenciamento de fenômenos sem
  necessariamente alegar identificação causal imediata.
\end{itemize}
\end{block}

\begin{block}{Legados Históricos e Mecanismos}
\protect\phantomsection\label{legados-histuxf3ricos-e-mecanismos}
\begin{itemize}
\tightlist
\item
  Argumentos de persistência histórica (Acemoglu et al.~2001; Nunn \&
  Wantchekon 2011; Homola et al.~2020) frequentemente sofrem de
  mecanismos de transmissão sub-especificados.
\item
  \textbf{Contribuição da Descrição}: Fornecer dados de resultados
  intermediários (\emph{intermediate outcome data}) e medir a
  ``meia-vida'' e atenuação dos legados (Cirone \& Pepinsky 2022).
\item
  \textbf{O Debate dos Länder Alemães} (Pepinsky et al.~2024): O
  controle descritivo da variação estadual em educação cívica desafia
  teses de persistência nazista direta.
\end{itemize}
\end{block}
\end{frame}

\section{Conclusão e Implicações
Metodológicas}\label{conclusuxe3o-e-implicauxe7uxf5es-metodoluxf3gicas}

\begin{frame}{Descrição como Parceira de Pleno Direito da Inferência
Causal}
\protect\phantomsection\label{descriuxe7uxe3o-como-parceira-de-pleno-direito-da-inferuxeancia-causal}
\begin{block}{Síntese dos Aportes do Paper}
\protect\phantomsection\label{suxedntese-dos-aportes-do-paper}
\begin{enumerate}
\tightlist
\item
  \textbf{Superação da Hierarquia Estigmatizante}: Descrição não é
  inferência causal de segunda classe.
\item
  \textbf{Rigor Conceitual}: Aplicação do critério de manipulabilidade
  (Holland 1986) e da escada de 6 degraus.
\item
  \textbf{Avaliação Tripartida}: Clareza, Comparabilidade e Completude
  como padrão de excelência.
\end{enumerate}

\begin{block}{Conclusão Epistêmica}
A boa descrição é uma arte científica teoricamente informada. Ela conecta teorias bem especificadas a métodos empíricos rigorosos, fundamentando e avaliando todo o empreendimento das ciências sociais.
\end{block}
\end{block}
\end{frame}

\begin{frame}{Guia de Avaliação para Pesquisadores e Teses}
\protect\phantomsection\label{guia-de-avaliauxe7uxe3o-para-pesquisadores-e-teses}
\begin{block}{Checklist para Leitura e Produção de Trabalhos
Descritivos}
\protect\phantomsection\label{checklist-para-leitura-e-produuxe7uxe3o-de-trabalhos-descritivos}
\begin{itemize}
\tightlist
\item[$\square$]
  \textbf{Teoria}: A pergunta descritiva deriva de uma teoria
  explicitamente declarada?
\item[$\square$]
  \textbf{Escada}: A análise está adequadamente posicionada na escada de
  perguntas (\(1 \to 6\))?
\item[$\square$]
  \textbf{Manipulabilidade}: O texto evita ambiguidades eufemísticas
  entre associação e causação?
\item[$\square$]
  \textbf{Contexto}: A mensuração considera o DGP local e a
  comparabilidade de significado?
\item[$\square$]
  \textbf{Método}: O método estatístico/qualitativo é adequado para
  descrever sem introduzir suposições contrafactuais indevidas?
\end{itemize}
\end{block}
\end{frame}




\end{document}
